\subsection{The complete physical return in the compact cubic block}

Fix the passive coefficient set and the margins in the raw sections.
Let \(z_*>0\) be the
unscaled height cap supplied by Corollary~\ref{rew:cor-high-derivative-signed};
its high-height estimate is used below after the cubic low-height analysis.

\begin{proposition}[Compact cubic coefficients]
\label{deg:prop-compact-cubic}
Fix \(U_0,G_0<\infty\), the original passive compact set and buffered
sections. For \(0\le U\le U_0\), \(|G|\le G_0\),
\(p=tU,\gamma=t^2G,s=t^3\), put \(\tau=\tau_{\rm cub}=t^2\).
Choose a positive unscaled raw cap \(0<h_c<z_*\) before the final
coefficient and multiplier bounds. On the entire original interval
\[
 I_\theta=\{x:a_\theta+x\in J_\theta^{\rm out},\quad
                 0<H(x)=z_\theta(\gamma;a_\theta+x)<h_c\},
 \qquad F(x)=x-m\Pi_{\theta,m}(x),
\]
the connected lower curvature receiver satisfies
\(\Pi''=0\Rightarrow\Pi'''<0\). Its bounded and growing descriptions
have actual resection overlap. The remaining admitted ghost raw tail has
\(0<\Pi'\le C(d_0)\tau^3\); its upper core and signed high complement
have bounded complete derivative and strictly overlap. After the common
multiplier shrink the same \(F'\) has at most two zeros, and \(F\)
has at most three isolated zeros, counted with multiplicity, on the whole
coefficient fibre. All original window, inverse-image, speed, stop and
height clips and intervening nonroot inputs are retained. Real-root and
fold endpoints are one-sided; the stopped fold continuation is excluded.
\end{proposition}

\begin{proof}
We construct normal sections on which the action controls curvature,
then continue their estimates to the original raw sections. Near the
roots, the two section variables must be estimated independently before
they are matched. The ghost tail and upper continuation require only a
bound on the complete first derivative. These estimates will cover one
connected physical interval to which the counting lemma applies.

In raw coordinates \(v-\gamma=\tau\xi,z=\tau^2\eta\), the exact field is
\begin{equation}\label{deg:compact-cubic-field}
 \dot\xi=-D,\quad\dot\eta=-\eta H,\qquad
 D=\tau P+Q\eta,\quad H=\xi B+C\tau^2(G+\xi)\eta,
\end{equation}
\begin{equation}\label{deg:compact-cubic-polynomials}
 P=1-U\xi+3G\xi^2+\xi^3,\quad B=1+\tau(2G+\xi),\quad
 Q=1-A\tau(G+\xi)+C\tau^2(G+\xi)^2.
\end{equation}
The unique simple negative root is \(-c\), by Descartes' rule, and
\begin{equation}\label{deg:compact-factorization}
 P(\xi)=(\xi+c)[(\xi-\alpha)^2-\Delta],\quad
 c-2\alpha=3G,\quad c(\alpha^2-\Delta)=1,\quad
 2\alpha c=U+c^{-1}.
\end{equation}
The four positive quantities \(c,c^{-1},\alpha,\alpha^{-1}\) are units.
A triple root would force \(P=(\xi+1)^3,U=-3\), outside this block.

Choose \(\ell_0\) so the original positive input window exceeds the
first-lip height at its raw port with a fixed margin. The actual lips
\(\eta=\tau\ell_0P\), with abscissae \(u,-r\), obey
\begin{equation}\label{deg:compact-lip-speed}
 \frac d{dt}\log\frac\eta P=-\xi B+\tau P'(1+Q\ell_0)
                              -C\tau^3(G+\xi)\ell_0P.
\end{equation}
Beyond a fixed positive threshold the signs are negative on the positive
lip and positive on the negative lip; their first hits are unique. At
\(a=\tau u\) the raw height is
\begin{equation}\label{deg:compact-lip-height}
 z=\ell_0[a^3+3G\tau a^2-U\tau^2a+\tau^3].
\end{equation}
Choose \(0<d_{\rm in}<d_0<\delta\), then \(\ell_0\delta\ll1\),
and finally \(\tau\ll d_{\rm in}\). Comparison with the organizing
arc and Lemma~\ref{lem:loc-capture} place captured low-height cycles at
\begin{equation}\label{deg:compact-raw-range}
                         0<a=\tau u<d_{\rm in}.
\end{equation}
The analytic carrier still includes all its original through inputs.
To continue them to the raw sections, set
\(v-\gamma=\sigma q\), \(\sigma=\pm1\), and
\(\mathcal W=z/q^2\). The graph equation becomes
\begin{equation}\label{deg:compact-raw-corridor}
 q\mathcal W_q=\mathcal W
 \frac{B_\sigma(q)+\sigma Cv_\sigma(q)q\mathcal W}
      {\ell(v_\sigma(q))\mathcal W+M(\sigma q)/q^2}-2\mathcal W.
\end{equation}
Here
\[
 v_\sigma(q)=\gamma+\sigma q,\qquad
 B_\sigma(q)=1+2\gamma+\sigma q,\qquad
 M(x)=x^3+3\gamma x^2-spx+s^2,
\]
and \(\ell\) is defined in \eqref{deg:central-coordinates}.
At fixed relative distance from a lip this is a small perturbation of
\(q\mathcal W_q=1-2\mathcal W\); constants around \(1/2\) give a
strict invariant strip to both original raw sections. The negative
continuation crosses its axis root at positive height. The original
windows and scalar order supply one through component.

\paragraph{Normal sections and independent root distances.}
In the bounded normal-section calculation below, \(p\) denotes the
positive axis footpoint; the fixed coefficient \(p=tU\) is represented
by \(U\) in \eqref{deg:compact-cubic-polynomials}.
Choose \(v>0\) with \(P\ge3/4\) on \([-4v,4v]\).
The limiting action exit for \(p=v\) lies in \([v/2,2v]\), because
both densities are \(s(1+O(v))\). Opposite strict action brackets and
Corollary~\ref{neg:cor-rootfree} give an actual comparison orbit,
independently of the return equation, with both footpoints in fixed
subintervals of \((0,3v)\).
Fix normal boxes and independent weak cuts \(p_w,r_w\) below these
footpoints. Scalar order selects the lower receiver, whose footpoints
have fixed positive lower bounds. The strong projections and normal
linearization give
\[
 \begin{gathered}
 \dot p=-\tau P(p),\quad\dot\zeta=-pB(p)\zeta,\quad\zeta_{\rm section}=u_+(p),\\
 J_+(p)=\int_{p_w}^p\frac{\xi B}{P}\,d\xi,\qquad
 J_-(r)=\int_{r_w}^r\frac{sB(-s)}{P(-s)}\,ds.
 \end{gathered}
\]
The projections have integrable kernels
\(C(1+T)^N e^{-(c_*-C\tau)T}\). To retain the higher rows, write
the attracting field as
\(V=(h-qy,-ny-ky^2)\), where \(h=O_{C^6}(\tau)\),
\(n\ge c_*>0\), and \(q,n,k\) have bounded finite derivatives.
Apply the variational induction in the proof of
Lemma~\ref{lem:normal-time-projection} to this attracting field.
The normal forcing always contains a height or a lower normal
variation, while pure tangential forcing has an \(O(\tau)\)
coefficient. Integrating the normal row first gives the same integrable
Volterra kernel. At phase order five it yields tangential rows
\(C(1+T)^Ne^{C\tau T}\) and normal rows
\(C(1+T)^Ne^{-(c_*-C\tau)T}\), using axis-flow order six.
A selected coefficient derivative requires mixed order seven of the
explicit coefficients and adds only polynomial time factors to this
majorant. Taylor division thus gives the required \(C^4\) units
\(u_\pm,u_\pm^{-1}\). The normal gap persists at the fold; these
coordinates are defined on the actual normal sections.

The action exponential must absorb derivative losses as either endpoint
approaches a root. We estimate those losses in five overlapping charts.
For \(\Delta=s_0^2\), \(a_0=\alpha-s_0\), \(X=a_0-p\), use
\((X,s_0/X)\) and \((s_0,X/s_0)\); the singular action integral is
\[
 \frac1{2s_0}\log\frac{\ell(X+2s_0)}{X(\ell+2s_0)},
                    \qquad s_0=0:\quad X^{-1}-\ell^{-1}.
\]
For \(\Delta=-\kappa^2\), use the left, central and right charts
\((X,\kappa/X)\), \(p=\alpha+\kappa z,|z|\le Z_0\), and
\((\widetilde X,\kappa/\widetilde X)\). Their arctangent action is
\(\kappa^{-1}[\arctan((p-\alpha)/\kappa)+\arctan(\ell/\kappa)]\).
Restoring the smooth positive density gives lower bounds
\(c/X-C\), \(cs_0^{-1}(1+|\log(X/s_0)|)-C\), or
\(c_{Z_0}/\kappa-C\). Every required ordinary phase derivative loses
only fixed powers of these independent variables. For example the real
inner loss is bounded by
\(s_0^{-N}z^{-N}e^{-c(1+|\log z|)/(2\tau s_0)}\), uniformly after
\(c/(2\tau_0s_{0,\max})>N+1\). Thus \(e^{-J_+/(2\tau)}\) pays
all root-distance losses, without comparing the two gaps.

\paragraph{Matching through the weak core.}
To compare the two actions, we need a core map defined all the way to
zero height. We construct that map on the whole chord below each
admitted orbit, before choosing an action cutoff. In raw height the
fixed weak-core equation is
\[
 \eta_\xi=f_\tau(\xi,\eta)
 =\frac{\eta(a+c_1\eta)}{T+Q\eta},\qquad
 a=\xi B,\quad c_1=C\tau^2(G+\xi),\quad T=\tau P.
\]
Here \(P\ge p_c>0\), \(Q\ge q_c>0\), and the coefficients are
bounded on a fixed enlarged weak slab. Direct differentiation gives
\[
 \partial_\eta f_\tau
 =\frac{aT+c_1\eta(2T+Q\eta)}{(T+Q\eta)^2},\qquad
 |\partial_\eta f_\tau|
 \le\frac{|a|}{T}+\frac{|c_1|}{Q}\le\frac C\tau,
 \qquad \eta\ge0.
\]
The constants use neither endpoint root distance.

Retain the actual intrinsic weak cuts from the projection construction
above. Write their raw parametrizations as
\(\iota_\pm(w)=(\psi_\pm(w),h_\pm(w))\), where \(w=\zeta_\pm\),
\(\psi_+(0)=p_w\), \(\psi_-(0)=-r_w\), and \(h_\pm(0)=0\).
The projection and normal linearization above, constructed by the
backward integrals of Lemma~\ref{lem:normal-time-projection} with the
displayed higher rows, give uniformly \(C^4\) inverse charts;
\(h_\pm'\) and the chart Jacobians
\(\Delta_\pm=\det D_{(v,\eta)}(\pi_\pm,\zeta_\pm)\),
in the respective radii \(v=\xi,-\xi\), are positive units.
Their axis values are one.
Use reversed physical time in the negative radial chart, so that both
attracting projected equations are
\(\dot\pi_\pm=-\tau P(\pm\pi_\pm)\).
The sections are \(\pi_+=p_w\), \(\pi_-=r_w\), while
\(\psi_-\) records the raw abscissa \(-v\).
Their exact projected axis equations give the transverse factors
\[
 \begin{aligned}
 A_+(w)&:=h_+'(w)-f_\tau(\iota_+(w))\psi_+'(w)
       =\frac{\tau P(p_w)}{D(\iota_+(w))\Delta_+(w)},\\
 A_-(w)&:=h_-'(w)-f_\tau(\iota_-(w))\psi_-'(w)
       =\frac{\tau P(-r_w)}{D(\iota_-(w))\Delta_-(w)}.
 \end{aligned}
\]
The reflection on the negative side reverses the abscissa and the
attracting time orientation, leaving both factors positive.
Since the weak buffers have fixed bounded height, \(c\tau\le A_\pm\le C\).

Fix any already admitted through orbit of the lower receiver, with weak
heights \(z>0\) and \(w>0\). Its core segment, the axis and the two
weak-section arcs from zero to these heights bound a strip. At either
cut the positive-time derivative of
\(\xi-\psi_\pm(h_\pm^{-1}(\eta))\) is \(-DA_\pm/h_\pm'<0\).
The strip lies to the left of the positive arc and to the right of the
negative arc, making these respectively incoming and outgoing.
Scalar uniqueness prevents a trajectory from crossing the
axis or this admitted orbit, and \(-\dot\xi\ge\tau p_c\) gives a
finite passage. Thus every point of \(0\le z_0\le z\) reaches the
negative arc inside this strip. The same actual intrinsic map
\(M_\tau\) is consequently defined on the entire chord, with
\(M_\tau(0)=0\) and \(0\le M_\tau(z_0)\le w\).
Uniqueness identifies these chord maps with each other and with the
local map of the weak-core flow.

Let \(V\) be the raw initial-height variation between the two actual
endpoint abscissae. Its equation \(V_\xi=(\partial_\eta f_\tau)V\)
on a slab of fixed length gives \(e^{-K/\tau}\le V\le e^{K/\tau}\).
Differentiating the actual endpoint equation yields
\[
 M_\tau'(z_0)=V\,
 \frac{A_+(z_0)}{A_-(M_\tau(z_0))}>0.
\]
Absorbing the fixed powers of \(\tau\) in the exponential constant,
\[
 e^{-K/\tau}\le M_\tau'(z_0)\le e^{K/\tau},\qquad
 \frac{M_\tau(z)}z=\int_0^1M_\tau'(tz)\,dt.
\]
The matching of these same intrinsic heights is
\(w_-=M_\tau(w_+)\), \(w_\pm=u_\pm e^{-J_\pm/\tau}\).
The positive \(C^4\) units \(u_\pm,u_\pm^{-1}\) therefore give
\[
 J_+-J_-
 =\tau\log\frac{M_\tau(w_+)}{w_+}
   +\tau\log\frac{u_+}{u_-},\qquad
 |J_+-J_-|\le K+C\tau\le K_0.
\]
Thus the actions differ by a bounded amount on every admitted through
orbit in the original independent section domains. The bound is obtained
before choosing an action cutoff and uses no return equation.

On the fixed weak core \(D\ge\tau p_c+q_c\eta\), the first four
normalized variational equations have bounds \(e^{K_j/\tau}\).
Factor the map just constructed at zero height as
\(M_\tau(z)=m_\tau zQ_\tau(z)\),
\(Q_\tau(0)=1,m_\tau^{\pm1}\le e^{K/\tau}\).
In particular, the exact quotient is
\[
 Q_\tau(z)=m_\tau^{-1}\int_0^1 M_\tau'(vz)\,dv.
\]
Four core variational equations bound its three ordinary derivatives.
Enlarge the exponential
constant to include this quotient, \(m_\tau^{-1}\), and the actual
input and output chord reserves; then choose \(J_*\) larger than
twice that constant with a fixed buffer. Thus
\(w_+=u_+(p)e^{-J_+/\tau}\) lies on the actual core chord and its
image lies in the negative weak image. Since
\(u_-(r)e^{-J_-/\tau}\downarrow0\) as \(r\uparrow c\), it has a
unique actual inverse. Taking logarithms on the product of separate
section domains gives
\begin{equation}\label{rc:strong-matching}
 \begin{gathered}
 J_+-J_-=C_\theta+\tau R(p,r),\qquad C_\theta=\tau\log m_\tau,\\
 R=\log u_+-\log u_-+\log Q_\tau(w_+),\\
 |C_\theta|\le K_0,\qquad\|R\|_{C^3_{p,r}}\le C(1+|\log\tau|).
 \end{gathered}
\end{equation}
Split the positive action exponential to absorb the core bounds and the
losses in the five root charts. Differentiation in the negative section
variable acts directly on \(u_-\); this avoids differentiating an
exponentially large inverse. Since the heights are intrinsic normal
coordinates, there is no moving initial \(\log P\) term.
On the actual matched graph, with \(P_+=P(p),P_-=P(-r)\),
\begin{equation}\label{rc:strong-variation}
 r_p=\frac{pB_+P_-}{rB_-P_+}(1+\rho),\qquad
 1+\rho=\frac{1-\tau P_+R_p/(pB_+)}{1+\tau P_-R_r/(rB_-)},
 \qquad\|\rho\|_{C^2_{p,r}}\le\epsilon_\tau:=C\tau(1+|\log\tau|).
\end{equation}
The polynomial multipliers and their derivatives are bounded; all
remaining denominators are units, including when \(P_+,P_-\) vanish
independently.

\paragraph{Curvature in the physical input.}
Let \(\Gamma_\pm(b)=a_\pm+c_\pm\phi_\pm(b)\), \(c_\pm>0\), be
these actual outward embeddings into the incoming raw height \(h\)
and divided output, with \(b=-p^2/2,S(b)=-r^2/2\).
The fixed normal-height equation
\(d(\eta-\xi^2/2)/d\log\eta=\tau\Phi\) has four bounded height
rows and clock \(O(1+|\log\tau|)\), independently of the axis gap.
For the exterior continuation retain \(\lambda_G=12\tau G/(1+2\tau G)\).
At each frozen bounded base \(q>0\),
\begin{equation}\label{deg:compact-section-integral}
 \int_1^{\delta/(\tau q)}\rho^{-(2-\lambda_G)}
 [\tau q^{-2}\rho^{-2}+\tau Uq^{-1}\rho^{-1}+\tau q\rho]
 (1+|\log\rho|^2)\frac{d\rho}\rho\le C\tau(1+q).
\end{equation}
The actual relative hierarchy yields
\(\phi_\pm=\mathrm{id}+O_{C^3}(\epsilon_\tau)\), including inverse
maps on their actual images. At an ordinary observer with exact small
gain \(a_\tau\asymp\tau^{2-O(\tau)}\), the normalized second and
third rows are \(a_\tau\chi''/\chi'_*,a_\tau^2\chi'''/\chi'_*\).
Only its ordinary \(C^3\) bounds are used. Outgoing and incoming
cocycles at every cut cancel before normalization; the output closing
retains \(f'(md)d'\) and its ordinary third chain rule.
Thus \(\widehat\Pi=\phi_-\circ S\circ\phi_+^{-1}\) is the actual
completed map up to its two retained positive affine gains.
The exact chain rule gives
\begin{equation}\label{deg:compact-finite-curvature}
 \phi_+'(b)\frac{P(p)}{p+r}\frac{\widehat\Pi''}{\widehat\Pi'}
 =3-\frac U{pr}+E(p,r),\qquad\|E\|_{C^1_{p,r}}\le\epsilon_\tau.
\end{equation}
To display the error, put \(W=p+r\),
\(K_{\rm mat}=B_+(1+\rho)/B_-\), and
\(L_\pm=\phi_\pm''/\phi_\pm'\), each in its own quadratic label.
Then
\[
 \begin{aligned}
 E={}&-\frac{(K_{\rm mat}-1)P_-'}{rW}
       -\frac{P_+}{pW}\partial_p\log K_{\rm mat}
       -\frac{K_{\rm mat}P_-}{rW}\partial_r\log K_{\rm mat}\\
    &+\frac{K_{\rm mat}P_-}{W}L_-
       -\frac{P_+}{W}L_+ .
 \end{aligned}
\]
Here \(P_-'\) uses the radial variable \(r\).
Since \(K_{\rm mat}-1=O_{C^2}(\epsilon_\tau)\) and
\(L_\pm=O_{C^1}(\epsilon_\tau)\), every term has the asserted
\(C^1\) bound on the product of the two normal-section domains:
\(P_\pm\) occur as multipliers and the remaining denominators are units.
Near a zero, \(U\ge U_*>0\), and restriction to the matched graph gives
\[
 \begin{aligned}
 \frac d{dp}\left(3-\frac U{pr}\right)
 &=\frac U{pr}\left(\frac1p+\frac{r_p}{r}\right)
 \ge c(1+r_p),\\
 \left|\frac d{dp}E(p,r(p))\right|
 &\le\epsilon_\tau(1+r_p).
 \end{aligned}
\]
Since \(w_p=(\phi_+(-p^2/2))_p<0\), every curvature zero is
simple and downward in the physical input.

On \(J_+\le J_{\rm rf}>J_*+1\),
both actions are bounded. The five singular integrals give a fixed
positive polynomial margin on the whole traversed segment and its
buffered exits. Corollary~\ref{neg:cor-rootfree}, with cubic coefficient
one, gives actual absolute \(C^3\) comparison there. The limiting map has
\(S_0'=P_-/P_+\),
\(S_0''/S_0'=(p+r)(3-U/(pr))/P_+\), and
\((pr)_b=-r/p-pP_-/(rP_+)<0\).
At \(U=3pr\) its \(S_0'''\) has a uniform negative margin; away
from that set the curvature has a nonzero margin. This transfers the
same crossing orientation to the bounded-action region. That comparison
strictly overlaps the high-action estimate for the same map.

On ghost fibres choose \(M\) above all roots and bottlenecks, large
for Lemma~\ref{deg:lem-weighted-cubic-tail}. Its actual growing
receiver \(u\ge2M,\tau u\le d_0\) has matched \(r_0\le r<c\)
and \eqref{deg:compact-curvature}, with error
\(C\tau(1+u)(1+|\log(\tau(1+u))|)^4\). Choose
\(M>2U_0/r_0\), then \(d_0\) with
\(Cd_0(1+|\log d_0|)^4<1/4\), then \(\tau\). Its curvature is
positive throughout this growing range.
To establish actual overlap, take a buffered band
\([2M-\beta,3M+\beta]\), \(0<\beta<M/2\), in the fixed normal
boxes. The lip projection is \(k_\tau(u)=u+O_{C^3}(\tau)\),
its linearized height is \(v_{\rm lip}=\tau\ell_0P(u)V_\tau(u)\),
with \(V_\tau\) a \(C^3\) unit. The same-orbit invariant gives
\begin{equation}\label{deg:compact-resection}
 \int_{k_\tau(u)}^{p_\tau(u)}\frac{qB(q)}{P(q)}\,dq
 =\tau\log\frac{u_+(p_\tau(u))}{v_{\rm lip}(u)},\qquad
 \|p_\tau-\mathrm{id}\|_{C^3}\le C\tau(1+|\log\tau|),\quad p_\tau'>0.
\end{equation}
The divergent bottleneck prefix cancels in this local integral. Its
\(p\)-derivative after subtracting \(\tau u_+'/u_+\) is a unit,
so the implicit equation gives all three rows. The closed band
\([2M+\beta,3M-\beta]\) has \(p_\tau<3M,u>2M\); the inverse
image of \(p=3M\) also has \(u>2M\), leaving no gap. These actual
orbits retain the original windows: their bounded normal labels have
raw gains \(\tau^{2-\lambda_G}\) times units, their reference images
lie inside the original passive organizing-arc buffers, and the raw
corridor continues them across the negative root at positive height.
Scalar order keeps the band in the comparison component. Resection
identifies orbits; each curvature estimate uses its own actual adapters.

\paragraph{The remaining raw ghost tail.}
On a ghost fibre, through trajectories can continue beyond the cutoff
\(\tau u=d_0\) of the growing estimate. The first-lip abscissa differs from the minimum
speed \(\sigma_\theta\); small central height may coexist with
\(a=\tau u>d_0\). Put \(a_b=d_0/2\), and let
\(x_*=\delta-\gamma\) be the actual positive raw port. For every
admitted \(a_b\le a\le x_*\), at fixed coefficients set
\[
 \begin{gathered}
 M_\tau=x^3+3\tau Gx^2-\tau^2Ux+\tau^3,\qquad B_0=1+2\tau G+x,\\
 Q_0=1-A(\tau G+x)+C(\tau G+x)^2,\quad
 \zeta_0=C(\tau G+x),\quad D_0=M_\tau+Q_0z.
 \end{gathered}
\]
The true graph is \(z_x=z(xB_0+\zeta_0z)/D_0\),
\(z(a)=\ell_0M_\tau(a)\). On \([a,x_*]\), \(M_\tau\ge c(d_0)>0\)
and \(Q_0\) is a positive unit. The graph satisfies
\(z_x\le Cx+Cz\), so its positive initial height and Gronwall
give \(z\le C\delta^2\), below the fixed upper chart boundary.
Its numerator is positive after the coefficient shrink, and hence
\(z\ge\ell_0M_\tau(a)>0\). This constructs the regular outward
passage, including its auxiliary zero-flight endpoint. Append the
observer only when the image lies in its original buffer; the auxiliary
endpoint needs no observer value when its image is outside that buffer.
Use \eqref{deg:scalar-polynomial-lip-slope}, with \(\sigma=\tau\),
\(K_{\rm lip}=K\), \(I_0=I\), and \(x_T=x_*\):
\begin{equation}\label{deg:compact-raw-tail-input}
 \Gamma_{+,a}'=-a\chi'(z_{\rm port})\frac{z_{\rm port}}{D_0(x_*,z_{\rm port})}
                   B_0(a)K_{\rm lip}(a)e^{I_0(a)}.
\end{equation}
The original lip transversality makes \(K_{\rm lip}\) a positive
unit. All admitted factors have bounds and reciprocals depending on
\(d_0\), so \(c(d_0)\le-\Gamma_{+,a}'\le C(d_0)\), or
\(c(d_0)\tau^2\le\Gamma_{+,b_+}'\le C(d_0)\tau^2\)
for \(b_+=-a^2/(2\tau^2)\).

We match this outward map to the negative branch by their actual central
heights. Differentiate the positive inward passage directly. For its height
\(\eta_0^+(u)\), \(W=\eta/(\tau P)\),
\(H_+=P'/(\xi B)\), \(\varpi=C\tau^3(G+\xi)P/(\xi B)\),
\begin{equation}\label{deg:compact-raw-tail-inward}
 \partial_u\log\eta_0^+=-\frac{uB_+}{\tau P_+}\Lambda_+,
 \quad\Lambda_+=\frac{K_{\rm lip}(a)}{1+Q_+\ell_0}
       \exp\int_u^0\frac{\xi B}{\tau P}
                         \frac{W(\varpi-Q)}{(1+QW)^2}\,d\xi.
\end{equation}
In the action clock \(dq=-\xi B\,d\xi/(\tau P)\),
\[
 W_q=W\left[-\frac{1+\varpi W}{1+QW}+\tau H_+\right],\qquad W(0)=\ell_0.
\]
On the original raw chart \(Q-\varpi>0\); the braces increase with
\(W\), and at \(\ell_0\) equal
\(-K_{\rm lip}/(1+Q\ell_0)\le-c_1\). Thus
\(W\le\ell_0e^{-c_1q}\) down to a fixed small scaled cut
\(\delta_c\), on the entire original raw range. The integral in
\eqref{deg:compact-raw-tail-inward} becomes the bounded integral
\(\int W(Q-\varpi)/(1+QW)^2\,dq\).
For the negative branch put \(P_-(v)=P(-v)=(c-v)L_-(v)\),
\(B_-(v)=B(-v)\), \(Q_-(v)=Q(-v)\), with
the exact unit \(L_-\) used in
Lemma~\ref{deg:lem-weighted-cubic-tail}. Its inverse inward graph is
\[
 \eta_v=f_-(v,\eta)
 =\frac{\eta\{vB_-+\zeta_-\eta\}}{\tau P_-+Q_-\eta},\qquad
 \zeta_-=-C\tau^2(G-v),\qquad
 \eta(r)=\tau\ell_0P_-(r).
\]
On \(\delta_c\le v\le r<c\), set
\(H_-=P_-'/ (vB_-)\) and
\(\varpi_-=\tau P_-\zeta_-/(vB_-)\).
The compact coefficient block gives bounded \(H_-\),
\(\varpi_-=O(\tau^3)\), and positive-unit \(Q_-\), even at
the simple root. The reflected ratio \(W_-\) has the same action-clock
equation and strict decay. Its complete moving-initial factor is
\[
 K_-=1-\tau H_-(1+\ell_0Q_-)+\ell_0\varpi_->c_K>0.
\]

Choose \(2\delta_c<r_0\). A fixed positive band below every
positive raw-tail label and the fixed negative band
\([\delta_c,r_0]\) supply action margins \(A_+,A_->0\).
Thus both inward graphs reach their central slab with heights
\(\eta_i\le C\tau e^{-c_1A_{\rm in}/\tau}\),
\(A_{\rm in}=\min(A_+,A_-)\). On that slab \(P_\pm\ge p_c>0\).
In the decreasing radial coordinate each graph satisfies
\(\eta_t\le C\tau\eta^2\); first exit keeps \(0<\eta\le2\eta_i\)
and gives an actual finite continuation to \(\xi=0\).
For \(f_+\) use the positive graph, with
\(\zeta_+=C\tau^2(G+v)\). The exact rational difference is
\[
 f_{\pm,\eta}-f_\pm/\eta
 =\frac{\eta\{\zeta_\pm\tau P_\pm-n_\pm Q_\pm\}}
             {(\tau P_\pm+Q_\pm\eta)^2},\qquad n_\pm=vB_\pm.
\]
Its central integral is bounded by \(C\eta_i/\tau^2\).
Together with the bounded noncentral ratio integral this proves the
positive \(\Lambda_+\) bound and, independently, the reflected identity
\begin{equation}\label{deg:compact-reflected-inward}
 \begin{gathered}
 -\partial_r\log\eta_0^-(r)=\frac{rB_-}{\tau P_-}\Lambda_-,\\
 \Lambda_-=
 \frac{K_-(r)}{1+\ell_0Q_-(r)}
       \exp\int_r^0(f_{-,\eta}-f_-/\eta)\,dv,\qquad
 0<c_\Lambda\le\Lambda_\pm,\Lambda_\pm^{-1}\le C_\Lambda.
 \end{gathered}
\end{equation}
No root factor has been inverted in these error bounds.

The ratio memory and the central linear cocycle give
\(\eta_0^-(r)=\tau\ell_0P_-(r)e^{-A_-(r)/\tau}E_-(r)\),
where \(E_-,E_-^{-1}\le C\) and
\(A_-(r)=\int_0^r vB_-/P_-\,dv\).
The simple root makes \(A_-(r)\to\infty\) as \(r\uparrow c\).
Hence the actual negative centre map is strictly decreasing with range
\begin{equation}\label{deg:compact-negative-centre-image}
              \eta_0^-([r_0,c))=(0,\eta_0^-(r_0)].
\end{equation}
Start with an admitted shot \(u_{\rm ref}\) in the fixed resection
band, whose negative match has \(r_{\rm ref}\ge r_0\).
For small \(\tau\), \(\tau u_{\rm ref}<a_b\).
The weighted inward construction, before either port observer is
attached, is strictly decreasing up to \(u_b=a_b/\tau\), since
\(\tau u\le a_b<d_0\). Thus
\(\eta_0^+(u_b)<\eta_0^-(r_{\rm ref})\), and the full negative
image gives a unique \(r_b\ge r_{\rm ref}\).
This auxiliary centre match does not require the formal \(a_b\) port
image to lie in an observer window.
Positive centre-height monotonicity and this full negative range then
give every original admitted \(a\ge a_b\) its unique negative match.
Scalar uniqueness at the same \(\xi=0\) identifies the joined graph
with the original through orbit; its negative output and closing
images therefore retain their original buffers. Matching them
gives \(S'=B_+P_-\Lambda_+/(B_-P_+\Lambda_-)\).
In the root collar the negative output embedding keeps its own endpoint
and exterior gains:
\[
 \begin{gathered}
 \Gamma_{-,b_-}'=d_c\theta_cX^{\beta_-}K_c(X,X^{\beta_-}),\quad X=c-r,\\
 d_c=\tau^{2-\lambda_G}U_c,\quad\theta_c=\tau^{\beta_-},\quad
 \beta_-=\tau P'(-c)/(cB(-c))>0.
 \end{gathered}
\]
Here \(U_c,K_c\) are bounded positive units. On the fixed complement
the ordinary negative exterior calculation gives
\(\Gamma_{-,b_-}'=\tau^{2-\lambda_G}\) times a bounded positive
unit. These estimates overlap on a fixed collar band and give
\(0<\Gamma_{-,b_-}'\le C\tau^2\) on the entire \(r_0\le r<c\),
since \(\tau^{-\lambda_G}\) is a unit. The closing contributes its actual
bounded positive first derivative. Now \(P_+(u)=\tau^{-3}M_\tau(a)
\ge c(d_0)\tau^{-3}\), while \(P_-\le C\). Therefore
\begin{equation}\label{deg:compact-raw-tail-derivative}
 0<\Pi'=\frac{\Gamma_{-,b_-}'}{\Gamma_{+,b_+}'}
           \frac{B_+P_-}{B_-P_+}\frac{\Lambda_+}{\Lambda_-}
                       \le C(d_0)\tau^3.
\end{equation}
The two absolute gains are retained in this ratio. The estimate covers
\(a_b\le a\le d_0\) and every original admitted \(a>d_0\), subject
to all the original window, speed and height restrictions. It therefore
continues the derivative bound beyond the growing estimate's cutoff.

\paragraph{The upper continuation.}
For the upper complement continue the fixed comparison outward to
\(\xi=\pm4v\). The fixed-normal fast invariants give positive-unit
heights there; scalar order preserves these lower bounds for all larger
inputs, and the field bounds them above through a fixed
\(\eta_{\max}>0\). With \(f=\eta H/D\), specialization of
\eqref{deg:scalar-raw-cocycle} gives
\begin{equation}\label{deg:compact-core-variation}
 f_\eta=\frac d{d\xi}\log\frac\eta D+
       \frac{\tau P'+(Q'+C\tau^2(G+\xi))\eta}{D}.
\end{equation}
Endpoint quotients are units, and \(D\ge\tau p_c+q_c\eta\) bounds
the integral even at exponentially small central height. The two outward
first variations have the same exact factor \(\tau^{2-\lambda_G}\);
their ratio and closing give \(0<\Pi'\le C_{\rm core}\).
The comparison orbit and a fixed action neighborhood lie in both the
lower and core regions. The comparison height \(\eta_{\rm comp}\)
is exponentially small.
The signed threshold is
\(H_H=\tau^3K_H(1+\tau U^2)\), \(\eta_H=H_H/\tau^2\).
After all fixed cuts and caps, choose \(\tau_0\) so
\(\eta_{\rm comp}<\eta_H<\eta_{\max}/2\) and
\(\tau^2\eta_{\max}<h_c/2\). Core and signed high then strictly
overlap on the original height axis.

\paragraph{Covering and counting the whole interval.}
The through lemma and increasing \(H,\sigma_\theta\) make
\(I_\theta\) connected. The comparison is admitted: its raw outward
images retain the original buffers, its central speed is
\(O(\tau^3+\tau^2\eta_{\rm comp})<d_{\rm out}\), and \(H<h_c\).
At its transverse positive lip every lower input has \(W<\ell_0\).
On a ghost fibre the fixed lower height of the original positive window
and the original lip reserve give \(W>\ell_0\) at the raw port.
Every input in \(I_\theta\) already has an actual through trajectory;
scalar order below the comparison gives \(W<\ell_0\) at its positive
lip. Thus a hit exists. Every hit has negative transverse derivative,
so it is unique. At the hit the radial derivative of \(W\) is positive,
whereas its input-height derivative is positive by scalar variation.
The implicit hit equation therefore gives \(a_x<0\) in the original
affine input, before any return-zero argument.
On real/fold fibres a through orbit crosses the incident inner root at
positive height, with \(W\to\infty\) on its inner side, giving the
same unique hit. An outer-component lip hit would force
\(\eta\le\tau\ell_0P\to0\) at the next root and a stopped passage.
Thus all lower through inputs have one of the actual descriptions above.
The bounded and growing ones have no gap by
\eqref{deg:compact-resection}; the growing one reaches \(a=d_0\),
and \eqref{deg:compact-raw-tail-derivative} continues every admitted tail.
The whole original carrier has the exact allocation
\begin{equation}\label{deg:compact-physical-cover}
 \begin{split}
 I_{\rm tail}&=\{x\in I_\theta:x\le x_{\rm comp},\ \Delta<0,\ a(x)>d_0\},\\
 I_{\rm lower}&=\{x\in I_\theta:x\le x_{\rm comp}\}\setminus I_{\rm tail},\\
 I_{\rm core}&=\{x\in I_\theta:x>x_{\rm comp},\ H(x)\le H_H\},\\
 I_{\rm high}&=\{x\in I_\theta:H(x)>H_H\},\qquad
 I_\theta=I_{\rm tail}\cup I_{\rm lower}\cup I_{\rm core}\cup I_{\rm high}.
 \end{split}
\end{equation}
They are intervals or empty clips, since \(a(x)\) decreases. Equality
\(a=d_0\) and the comparison belong to the lower piece; \(H=H_H\)
belongs to the core. Closed action dividers are retained, and all estimate
packages have the strict overlaps above. Natural stopped endpoints and
the infinite-time fold branch are excluded.
Choose \(m\) last with
\(m\max(C_{\rm core},C_H,C(d_0)\tau_0^3)<1/2\).
Then \(F'>1/2\) on tail, core and high. On the lower receiver the
exact two-gain composition gives
\[
 F=a_+-a_\theta-ma_-+c_+w-mc_-\widehat\Pi(w),\qquad
 F''=-\frac{mc_-}{c_+^2}\widehat\Pi'',\quad
 F'''=-\frac{mc_-}{c_+^3}\widehat\Pi'''.
\]
With \(J=I_{\rm lower}\), Lemma~\ref{sharp:upper-curvature-count}
applies on the connected \(I_\theta\): the positive factors preserve
the lower downward curvature crossings, and \(F'>1/2\) on
\(I_\theta\setminus J\). It gives at most one curvature zero on \(J\),
at most two zeros of \(F'\), counted with multiplicity, and at most
three isolated zeros of the
same physical \(F\) on \(I_\theta\), counted with multiplicity.
All choices depend only on the fixed coefficient block and section
margins. The estimates allow the two gains and the two root gaps to
vary independently.
\end{proof}

